\documentclass[12pt]{article}
\usepackage{styles/cwilliams-standard}

\setclass{\MLONE}
\settitle{Homework 9}

\begin{document}

\maketitlepage

\section{Estimating Probabilities (1)}

\begin{figure}[H]
    \centering
    \includegraphics[width=0.9\textwidth]{images/ML9_q1.jpeg}
\end{figure}

\section{Naïve Beyes Approach (1)}

\begin{figure}[H]
    \centering
    \includegraphics[width=0.85\textwidth]{images/ML9_q2.jpeg}
\end{figure}

\section{Estimating Probabilities w/ Laplace Smoothing}

\begin{figure}[H]
    \centering
    \includegraphics[width=0.75\textwidth]{images/ML9_q3.jpeg}
\end{figure}

\section{Naïve Beyes Approach (2)}

\begin{figure}[H]
    \centering
    \includegraphics[width=0.9\textwidth]{images/ML9_q4.jpg}
\end{figure}

\section{Comparison of 2 and 4}

Estimating probabilities with Laplace Smoothing is better. Laplace smoothing prevents a `0 term' from destroying or completely altering
the outcome like in Question 2, when estimating the posterior of $P(- | vals) \propto \frac{3}{5} \cdot \frac{2}{5} \cdot \frac{0}{5} \cdot \frac{5}{10} = 0$. That
single 0 term made the model predict with 100\% that the sample was `+' when in reality it should have been `-'.

With Laplace Smoothing, $P(- | vals) \propto \frac{3}{9} \cdot \frac{2}{9} \cdot \frac{0}{9} \cdot \frac{5}{10} = 0.526$, which allows the model
to properly compare both classes. In general, Laplace smoothing allows unseen feature-class combinations don't receive zero probability.

\section{Estimating Probabilities (2)}

\begin{figure}[H]
    \centering
    \includegraphics[width=0.85\textwidth]{images/ML9_q6.jpeg}
\end{figure}

\section{Naïve Beyes Approach (3)}

\begin{figure}[H]
    \centering
    \includegraphics[width=0.9\textwidth]{images/ML9_q7.jpeg}
\end{figure}

\section{Bayesian Network}

\begin{figure}[H]
    \centering
    \includegraphics[width=0.9\textwidth]{images/ML9_q8_1.jpeg}
\end{figure}

\begin{figure}[H]
    \centering
    \includegraphics[width=1\textwidth]{images/ML9_q8_2.jpeg}
\end{figure}

\end{document}